\section{The unbalanced fixed-determinant ladders}
\label{sec:tpc45-determinant-ladders}

The upper endpoint \(N_-/\ell_+\) in
\eqref{eq:tpc45-balanced-band} is the natural boundary of the fiber
argument.  We now retain only the uniqueness condition
\(q>\sqrt{N_+}\) and allow the cofactor \(r=n/q\) to be smaller than the
source-prime scale.

Fix an orbit time \(j\), a residual label \(s\), and a divisor
\(d\mid s\).  Put \(r=s/d\).  Every row that deletes to \(s\) solves
\begin{equation}
 rq-dj\ell=h.
 \label{eq:tpc45-unbalanced-determinant}
\end{equation}
Let \(g=(r,dj)\).  If \(g\nmid h\), there are no such rows.  Otherwise,
for any one integral solution \((q_0,\ell_0)\), the complete ladder is
\begin{equation}
 \boxed{
 q=q_0+\frac{dj}{g}t,
 \qquad
 \ell=\ell_0+\frac r g t,
 \qquad t\in\Z.}
 \label{eq:tpc45-full-determinant-ladder}
\end{equation}
The two affine forms have fixed nonzero cross determinant
\begin{equation}
 \frac{dj}{g}\ell_0-\frac r g q_0=-\frac hg.
 \label{eq:tpc45-ladder-cross-determinant}
\end{equation}
In the physical packet both \(q\) and \(\ell\) must be prime, in their
respective support windows.  This primality condition deletes points
from \eqref{eq:tpc45-full-determinant-ladder}; it does not shorten the
ambient ladder available to a coefficient-blind argument.

\begin{proposition}[Dyadic residual-fiber capacity]
\label{prop:tpc45-unbalanced-fiber-capacity}
Let \(\cR\subset[R,2R]\) be a dyadic cofactor block.  For fixed
\((j,s)\), the number of rows with \(r=s/d\in\cR\),
\(q>\sqrt{N_+}\), and descended label \(s\) is at most
\begin{equation}
 C_h
 \sum_{\substack{d\mid s\\s/d\in\cR}}
 \left(1+\frac{L_0d}{s}\right)
 \le C_h\tau(s)\left(1+\frac{L_0}{R}\right).
 \label{eq:tpc45-unbalanced-fiber-bound}
\end{equation}
\end{proposition}

\begin{proof}
For fixed \(d\), the step of the \(\ell\)-form in
\eqref{eq:tpc45-full-determinant-ladder} is \(r/g\).  Since
\(g\mid h\) and the source-prime window has length \(O(L_0)\), the
number of integral parameters is
\[
 O\!\left(1+\frac{gL_0}{r}\right)
 =O_h\!\left(1+\frac{L_0d}{s}\right).
\]
Sum over the eligible divisors \(d\mid s\), and use \(r\ge R\).
\end{proof}

At the original equality endpoint,
\begin{equation}
 L_0=X^{99979/210000+o(1)},
 \qquad
 J=X^{133/400+o(1)}.
 \label{eq:tpc45-L-J-scales}
\end{equation}
At the old threshold \(q\asymp Q\), where the cofactor scale is
\(R\asymp J\), the unsieved ladder capacity is already
\begin{equation}
 \frac{L_0}{J}
 =X^{15077/105000+o(1)}.
 \label{eq:tpc45-best-unbalanced-capacity}
\end{equation}
It exceeds the endpoint allowance \(K_B=X^{1/400+o(1)}\) by the fixed
power
\begin{equation}
 \frac{15077}{105000}-\frac1{400}
 =\frac{29629}{210000}>0.
 \label{eq:tpc45-best-capacity-deficit}
\end{equation}
For \(r=O(1)\), the ambient ladder has length
\(X^{99979/210000+o(1)}\).

The parametrization also explains why a standard upper-bound sieve is
not the missing theorem.  It may remove logarithmic powers from the
number of parameters for which both affine forms are prime, but the
power-sized factor \(L_0/R\) remains.  More importantly, the physical
coefficients on the surviving parameters are neither identical to an
unweighted prime-pair count nor controlled by their cardinality.  The
needed input is a weighted Bessel or dispersion estimate for the actual
source and mask coefficients along
\eqref{eq:tpc45-full-determinant-ladder}.

Equation \eqref{eq:tpc45-unbalanced-determinant} is a literal
fixed-determinant equation, not a heuristic analogy.  Nevertheless, a
theorem for two arbitrary coefficient sequences and otherwise smooth
weights cannot be invoked after silently inserting, in addition, the
prime restrictions on \(q,\ell\), the M\"obius cofactor weight, and the
moving source/mask coefficient.  Those quantifiers form the next
arithmetic interface.
